\section{Experimental setup}

\subsection{Benchmark design and reference solutions}

The main benchmark contains \MainInstanceCount{} combinations of a precedence
graph and a prescribed workstation count, drawn from the
\MainFamilyCount{} classical SALBP families in
Table~\ref{tab:benchmark-inventory}. Processing times and precedence arcs come
from the SALBP instance files; task powers are the corresponding vectors used
in the power-peak SALBP benchmark. Each case is evaluated under
\(\Qmax^{\mathrm{top}}\) and \(\Qmax^{\mathrm{avg}}\), computed from the task
powers by Equations~\eqref{eq:qmax-one} and \eqref{eq:qmax-two}.

% This table is generated by analysis/analyze_manuscript_results.py.
% Do not edit numerical values manually.
\begin{table}[t!]
\centering
\caption{Benchmark set. Each case combines a precedence graph with a prescribed number of workstations $m$.}
\label{tab:benchmark-inventory}
\footnotesize
\renewcommand{\arraystretch}{1.15}
\setlength{\tabcolsep}{6pt}
\begin{tabular}{@{} l c c c c @{\quad} l c c c c @{}}
\toprule
Graph & $n$ & Cases & $m_{\min}$ & $m_{\max}$ & Graph & $n$ & Cases & $m_{\min}$ & $m_{\max}$ \\
\midrule
BOWMAN & 8 & 1 & 5 & 5 & MANSOOR & 11 & 3 & 2 & 4 \\
BUXEY & 29 & 6 & 7 & 13 & MERTENS & 7 & 4 & 2 & 6 \\
GUNTHER & 35 & 6 & 7 & 14 & MITCHELL & 21 & 3 & 3 & 8 \\
HESKIA & 28 & 4 & 3 & 8 & ROSZIEG & 25 & 4 & 4 & 10 \\
JACKSON & 11 & 5 & 3 & 8 & SAWYER & 30 & 8 & 5 & 14 \\
JAESCHKE & 9 & 4 & 3 & 8 & WARNECKE & 58 & 13 & 14 & 31 \\
LUTZ2 & 89 & 11 & 24 & 49 &  &  &  &  &  \\
\addlinespace[3pt]
\multicolumn{10}{@{}c}{\textit{Total: \MainFamilyCount{} precedence graphs and \MainInstanceCount{} benchmark cases}} \\
\bottomrule
\end{tabular}
\end{table}


The cycle-time optimum without a power cap is required to measure the cost of
the new constraint. We use the values in the standard Scholl SALBP benchmark
\citep{SchollBenchmark1993}, whose
SALBP-2 sheet reports numeric \(c^*\) values and marks unknown
cases with an interval \([\mathrm{LB},\mathrm{UB}]\). All 35 SALBP-2 source
values in the No-Peak CSV are numeric \(c^*\) values and match \texttt{Best}.
The remaining 37 cases start from a feasible cycle time in the SALBP-1 sheet;
they are treated as upper bounds only until the No-Peak SAT run proves the
same or a smaller cycle time. The two HESKIA timeout rows are closed by the
workload lower bound \(\lceil\sum_i t_i/m\rceil\), which equals the feasible
cycle time in both cases. Thus all 72 denominators are independently
certified SALBP-2 optima without using the local Direct-BBR reimplementation.

% This table is generated by analysis/analyze_manuscript_results.py.
% Do not edit numerical values manually.
\begin{table}[t!]
\centering
\caption{Independent provenance of the uncapped SALBP-2 reference values.}
\label{tab:no-peak-reference}
\footnotesize
\renewcommand{\arraystretch}{1.15}
\setlength{\tabcolsep}{5pt}
\begin{threeparttable}
\begin{tabular}{@{} l c l @{}}
\toprule
Reference source & Cases & Certification role \\
\midrule
Scholl SALBP-2 benchmark ($c^*$) & 35 & Published exact SALBP-2 optimum \\
Scholl SALBP-1 feasible $c$ + No-Peak SAT & 35 & SAT reports optimality \\
Scholl SALBP-1 feasible $c$ + workload LB & 2 & $\lceil\sum_i t_i/m\rceil = c$ \\
\bottomrule
\end{tabular}
\begin{tablenotes}[flushleft]
\scriptsize
\item The SALBP-1 column is a feasible upper bound for SALBP-2, not an optimum certificate by itself. Benchmark values are from \citet{SchollBenchmark1993}.
\end{tablenotes}
\end{threeparttable}
\end{table}


The No-Peak SAT campaign reports optimality for 63 cases within its recorded
runs. The seven WARNECKE timeout rows are already exact values published in
the SALBP-2 benchmark, and the two HESKIA timeout rows have the workload
lower-bound certificate described above. The BBR paper remains relevant
related work, but no locally modified BBR executable or runtime is used as a
baseline or as the source of the uncapped denominator.

\subsection{Methods and computational protocol}

The locked two-phase SAT campaign compares three encodings on the direct
precedence arc set \(E\): the baseline non-overlap encoding, the same-station
merger (SM), and the task-pair timing variant (SM-T). A separate archived
campaign compares \(E\) with its transitive closure \(E^*\); it is used only
for the edge-set ablation. The solver comparison uses Gurobi-MIP, CPLEX-MIP
and CPLEX-CP as exact baselines. The representative SAT encoding is selected
within each cap by proof coverage, with proof time used only to break a tie.
Additional source-labelled variants are retained in the archive but do not
support the main performance claims.

Every capped run has a 3,600-second wall-clock cutoff. SAT uses one thread;
the commercial solvers use their default automatic thread policy, as in the
recorded experiment configuration. The two-phase SAT search uses 50,000
conflicts per Phase-I feasibility probe and unrestricted calls in Phase~II.
Gurobi uses a 4\,GB soft memory limit. A run is counted as \emph{proven} only
when the solver reports \textsc{optimal}; a numeric value returned at timeout
or out-of-memory termination is treated as a solver-reported feasible upper
bound.
% TODO(author): insert the archived processor, operating-system and exact
% solver-version metadata before submission.

\subsection{Measures and research questions}

For each case and cap, the best known solution (BKS) is the smallest reported
feasible cycle time. It is labelled \emph{OPT} only if at least one method
proves the same value and \emph{UB} otherwise. For method \(a\), solution
quality is measured by
\[
  \operatorname{RPD}_a =
  100\,\frac{C_a-C_{\mathrm{BKS}}}{C_{\mathrm{BKS}}}.
\]
Runtime summaries report the median and geometric mean over runs proving
optimality. PAR-2 assigns 7,200 seconds to a run that does not prove
optimality. Pairwise speed ratios are computed only when both methods prove
the same optimum. The effect of a cap is
\[
  \Delta_Q =
  100\,\frac{C_{\mathrm{BKS}}(Q)-C^*_{\mathrm{NoPeak}}}
  {C^*_{\mathrm{NoPeak}}},
\]
where \(C^*_{\mathrm{NoPeak}}\) is the independently certified uncapped
SALBP-2 optimum. If a capped BKS is not proven, \(\Delta_Q\) describes its
best known feasible increase rather than an exact optimality loss.

The analysis addresses four questions: how strongly each cap changes cycle
time (RQ1); how the SAT encodings, two-phase search trajectories and
precedence arc sets affect solution and proof performance (RQ2); how the
selected SAT encodings compare with MIP and CP (RQ3); and how proof coverage
varies across basic instance characteristics (RQ4). All numerical tables and
figure coordinates below are generated from the five primary source CSV files by
\texttt{analysis/analyze\_manuscript\_results.py}.

\section{Results and discussion}

\subsection{RQ1: cycle-time effect of the power caps}

% This table is generated by analysis/analyze_manuscript_results.py.
% Do not edit numerical values manually.
\begin{table}[t!]
\centering
\caption{Best-known cycle time increase under each peak power threshold relative to the independently certified uncapped SALBP-2 optimum.}
\label{tab:cap-impact}
\footnotesize
\renewcommand{\arraystretch}{1.15}
\setlength{\tabcolsep}{3.5pt}
\begin{threeparttable}
\begin{tabular*}{\textwidth}{@{\extracolsep{\fill}} l c c c c c c c @{}}
\toprule
Threshold & Feasible BKS & BKS proven & Same $C$ & Higher $C$ & Mean $\Delta C$ & Mean $\Delta C$ (\%) & Median $\Delta C$ (\%) \\
\midrule
$Q^{\mathrm{top}}_{\max}$ & 72 & 45 & 6 & 66 & 7.07 & 16.49\% & 14.29\% \\
$Q^{\mathrm{avg}}_{\max}$ & 72 & 25 & 0 & 72 & 30.32 & 60.99\% & 62.20\% \\
\bottomrule
\end{tabular*}
\begin{tablenotes}[flushleft]
\scriptsize
\item Feasible BKS counts cases with a solver-reported feasible solution. BKS proven counts cases in which at least one solver proves that value optimal; the remaining capped BKS values are valid feasible upper bounds. The uncapped denominator is certified independently from the Scholl SALBP-2 values, No-Peak SAT proofs, and two workload lower-bound closures, so $\Delta C$ is an exact optimality increase only when the capped BKS is also proven.
\end{tablenotes}
\end{threeparttable}
\end{table}


The top-\(m\)-based cap admits a feasible BKS for every case; 45 are proven
optimal. Six cases retain their uncapped optimum, whereas 66 require a larger
cycle time. The mean and median increases are 16.49\% and 14.29\%,
respectively. The average-based cap is uniformly tighter in this benchmark.
Its feasible BKS is larger than the uncapped optimum in all 72 cases, with a
mean increase of 60.99\% and a median of 62.20\%; 25 of these BKS values are
proven optimal.

% This figure is generated by analysis/analyze_manuscript_results.py.
% Do not edit numerical coordinates manually.
\begin{figure}[H]
\centering
\begin{minipage}{0.49\textwidth}
\centering
\begin{tikzpicture}
\begin{axis}[width=\linewidth,height=57mm,xmin=0,xmax=95,ymin=0,ymax=95,xtick={0,20,40,60,80},ytick={0,20,40,60,80},xlabel={Increase under $Q^{\mathrm{top}}_{\max}$ (\%)},ylabel={Increase under $Q^{\mathrm{avg}}_{\max}$ (\%)},grid=major,grid style={gray!22,densely dotted},tick label style={font=\scriptsize},label style={font=\small}]
\addplot[domain=0:95,samples=2,no marks,gray!60,densely dashed] {x};
\addplot[only marks,mark=*,mark size=1.6pt,mark options={draw=PlotBlue!90!black,fill=PlotBlue,fill opacity=0.6,draw opacity=0.95}] coordinates {(35.294118,47.058824) (5.882353,55.882353) (3.125000,53.125000) (14.285714,60.714286) (11.111111,55.555556) (2.127660,51.063830) (0.000000,53.658537) (6.250000,68.750000) (9.090909,70.454545) (10.000000,62.500000) (4.166667,62.500000) (6.349206,66.666667) (11.111111,77.777778) (1.169591,56.432749) (5.859375,51.562500) (10.243902,58.536585) (23.255814,73.643411) (6.250000,31.250000) (8.333333,41.666667) (10.000000,40.000000) (11.111111,33.333333) (42.857143,42.857143) (7.692308,30.769231) (10.000000,30.000000) (0.000000,25.000000) (33.333333,33.333333) (14.285714,76.190476) (15.000000,80.000000) (15.789474,84.210526) (16.666667,77.777778) (23.529412,82.352941) (25.000000,87.500000) (20.000000,80.000000) (21.428571,78.571429) (23.076923,76.923077) (25.000000,75.000000) (27.272727,72.727273) (0.000000,13.978495) (1.612903,22.580645) (0.000000,18.750000) (0.000000,40.000000) (10.000000,40.000000) (28.571429,28.571429) (50.000000,50.000000) (5.714286,48.571429) (14.285714,61.904762) (14.285714,64.285714) (7.142857,57.142857) (3.125000,43.750000) (4.761905,61.904762) (12.500000,68.750000) (11.764706,58.823529) (16.129032,61.290323) (21.428571,67.857143) (26.923077,69.230769) (24.000000,64.000000) (0.000000,41.538462) (2.127660,51.063830) (7.317073,58.536585) (24.324324,82.882883) (25.000000,82.692308) (27.173913,84.782609) (27.380952,82.142857) (30.379747,84.810127) (30.263158,84.210526) (30.136986,83.561644) (33.333333,85.507246) (34.848485,86.363636) (34.375000,85.937500) (35.000000,83.333333) (37.500000,83.928571) (39.622642,83.018868)};
\end{axis}
\end{tikzpicture}
\par\smallskip\footnotesize (a) Within-case comparison
\end{minipage}
\hfill
\begin{minipage}{0.49\textwidth}
\centering
\begin{tikzpicture}
\begin{axis}[width=\linewidth,height=57mm,xmin=0.48,xmax=0.92,ymin=0,ymax=95,xtick={0.5,0.6,0.7,0.8,0.9},ytick={0,20,40,60,80},xlabel={Threshold tightness $s_Q$},ylabel={BKS increase (\%)},yticklabel style={font=\scriptsize},xticklabel style={font=\scriptsize},label style={font=\small},grid=major,grid style={gray!22,densely dotted},legend style={font=\scriptsize,at={(0.02,0.98)},anchor=north west,draw=gray!40,fill=white,fill opacity=0.95}]
\addplot[only marks,mark=*,mark size=1.6pt,mark options={draw=PlotBlue!90!black,fill=PlotBlue,fill opacity=0.7}] coordinates {(0.500000,35.294118) (0.500000,5.882353) (0.500000,3.125000) (0.501241,14.285714) (0.500000,11.111111) (0.500000,2.127660) (0.500000,0.000000) (0.501155,6.250000) (0.500000,9.090909) (0.500914,10.000000) (0.500000,4.166667) (0.501618,6.349206) (0.501425,11.111111) (0.505376,1.169591) (0.500000,5.859375) (0.502732,10.243902) (0.501650,23.255814) (0.507246,6.250000) (0.505051,8.333333) (0.500000,10.000000) (0.500000,11.111111) (0.500000,42.857143) (0.500000,7.692308) (0.500000,10.000000) (0.500000,0.000000) (0.503067,33.333333) (0.500000,14.285714) (0.500000,15.000000) (0.500000,15.789474) (0.500000,16.666667) (0.500000,23.529412) (0.500412,25.000000) (0.500000,20.000000) (0.500000,21.428571) (0.500000,23.076923) (0.500000,25.000000) (0.500000,27.272727) (0.500000,0.000000) (0.500000,1.612903) (0.503546,0.000000) (0.512195,0.000000) (0.500000,10.000000) (0.500000,28.571429) (0.503497,50.000000) (0.505376,5.714286) (0.502857,14.285714) (0.501742,14.285714) (0.500000,7.142857) (0.503650,3.125000) (0.500000,4.761905) (0.500000,12.500000) (0.500000,11.764706) (0.500000,16.129032) (0.500000,21.428571) (0.501529,26.923077) (0.500000,24.000000) (0.500000,0.000000) (0.500000,2.127660) (0.502326,7.317073) (0.500876,24.324324) (0.500000,25.000000) (0.500730,27.173913) (0.500659,27.380952) (0.500000,30.379747) (0.500000,30.263158) (0.500000,30.136986) (0.500554,33.333333) (0.500000,34.848485) (0.500000,34.375000) (0.500000,35.000000) (0.500000,37.500000) (0.500000,39.622642)};
\addlegendentry{$Q^{\mathrm{top}}_{\max}$}
\addplot[only marks,mark=triangle*,mark size=1.9pt,mark options={draw=PlotOrange!90!black,fill=PlotOrange,fill opacity=0.7}] coordinates {(0.614458,47.058824) (0.718391,55.882353) (0.704787,53.125000) (0.694789,60.714286) (0.683721,55.555556) (0.759843,51.063830) (0.743056,53.658537) (0.729792,68.750000) (0.725738,70.454545) (0.714808,62.500000) (0.751880,62.500000) (0.747573,66.666667) (0.740741,77.777778) (0.795699,56.432749) (0.760870,51.562500) (0.737705,58.536585) (0.699670,73.643411) (0.768116,31.250000) (0.717172,41.666667) (0.674603,40.000000) (0.646667,33.333333) (0.593750,42.857143) (0.791667,30.769231) (0.734694,30.000000) (0.654930,25.000000) (0.558282,33.333333) (0.700617,76.190476) (0.698413,80.000000) (0.696360,84.210526) (0.691203,77.777778) (0.689024,82.352941) (0.684774,87.500000) (0.677591,80.000000) (0.670940,78.571429) (0.663539,76.923077) (0.654182,75.000000) (0.641204,72.727273) (0.875000,13.978495) (0.781250,22.580645) (0.737589,18.750000) (0.902439,40.000000) (0.768293,40.000000) (0.619048,28.571429) (0.559441,50.000000) (0.817204,48.571429) (0.748571,61.904762) (0.707317,64.285714) (0.706215,57.142857) (0.810219,43.750000) (0.759259,61.904762) (0.731034,68.750000) (0.702290,58.823529) (0.690141,61.290323) (0.679739,67.857143) (0.669725,69.230769) (0.658046,64.000000) (0.808824,41.538462) (0.752632,51.063830) (0.734884,58.536585) (0.688266,82.882883) (0.685246,82.692308) (0.677372,84.782609) (0.670619,82.142857) (0.667085,84.810127) (0.664663,84.210526) (0.661290,83.561644) (0.658915,85.507246) (0.655650,86.363636) (0.653292,85.937500) (0.647399,83.333333) (0.640255,83.928571) (0.633218,83.018868)};
\addlegendentry{$Q^{\mathrm{avg}}_{\max}$}
\end{axis}
\end{tikzpicture}
\par\smallskip\footnotesize (b) Relative increase versus $s_Q$
\end{minipage}
\caption{Best-known cycle time increases under the two peak power threshold constructions. Panel~(a) compares the 72 benchmark cases with a feasible BKS under both thresholds; the dashed line denotes equality. Panel~(b) shows all 144 benchmark-case/threshold records.}
\label{fig:cap-effect}
\end{figure}


Figure~\ref{fig:cap-effect}(a) makes the within-case contrast explicit:
\(\Qmax^{\mathrm{avg}}<\Qmax^{\mathrm{top}}\) in all 72 cases, and its BKS
cycle time is higher in 68 cases and equal in four. Panel~(b) shows that
\(\Qmax^{\mathrm{top}}\) is constructed near the midpoint of its lower and
upper power references, while the tightness of \(\Qmax^{\mathrm{avg}}\)
varies much more across instances. The large separation between the two
series is therefore an effect of the threshold construction rather than a
change in the underlying SALBP data.

The corrected SAWYER case with \(n=30\) and \(m=12\) illustrates both the
magnitude and the status convention. Its initial probe horizon is \(C_0=54\)
and its exact No-Peak optimum is 28. The BKS values are 34 under
\(\Qmax^{\mathrm{top}}\) and 47 under \(\Qmax^{\mathrm{avg}}\), corresponding
to increases of 21.43\% and 67.86\%. All methods that reach these two values
terminate without a proof, so both are reported as feasible UBs.

\subsection{RQ2: SAT encoding and two-phase search}

% This table is generated by analysis/analyze_manuscript_results.py.
% Do not edit numerical values manually.
\begin{table}[t!]
\centering
\caption{End-of-run performance of all six Direct and Two-phase SAT configurations on $E$.}
\label{tab:sat-main-results}
\footnotesize
\renewcommand{\arraystretch}{1.15}
\setlength{\tabcolsep}{5pt}
\begin{threeparttable}
\begin{tabular}{@{} l c c c c @{}}
\toprule
Configuration & BKS & Proven & Mean RPD & PAR-2 (s) \\
\midrule
\multicolumn{5}{@{}l}{\textit{Top-$m$-based threshold ($Q^{\mathrm{top}}_{\max}$)}} \\
\addlinespace[2pt]
D-CSE/$E$ & 70 / 72 & \textbf{44 / 72} & 0.02\% & 2919.2 \\
2P-CSE/$E$ & 69 / 72 & \textbf{44 / 72} & 0.09\% & 2933.5 \\
\addlinespace[1pt]
D-SM/$E$ & \textbf{71 / 72} & \textbf{44 / 72} & \textbf{0.01\%} & \textbf{2911.2} \\
2P-SM/$E$ & 70 / 72 & 43 / 72 & 0.02\% & 2965.5 \\
\addlinespace[1pt]
D-SM-T/$E$ & 70 / 72 & \textbf{44 / 72} & 0.08\% & 2918.7 \\
2P-SM-T/$E$ & \textbf{71 / 72} & 43 / 72 & \textbf{0.01\%} & 2988.5 \\
\addlinespace[4pt]
\multicolumn{5}{@{}l}{\textit{Average-based threshold ($Q^{\mathrm{avg}}_{\max}$)}} \\
\addlinespace[2pt]
D-CSE/$E$ & 62 / 72 & 22 / 72 & 0.25\% & 5008.4 \\
2P-CSE/$E$ & 61 / 72 & 23 / 72 & 0.20\% & 4957.7 \\
\addlinespace[1pt]
D-SM/$E$ & \textbf{64 / 72} & 23 / 72 & 0.18\% & 4954.6 \\
2P-SM/$E$ & 62 / 72 & 23 / 72 & 0.14\% & 4956.9 \\
\addlinespace[1pt]
D-SM-T/$E$ & 62 / 72 & 23 / 72 & 0.25\% & 4951.5 \\
2P-SM-T/$E$ & \textbf{64 / 72} & \textbf{24 / 72} & \textbf{0.11\%} & \textbf{4903.1} \\
\bottomrule
\end{tabular}
\begin{tablenotes}[flushleft]
\scriptsize
\item Configuration IDs follow Table~\ref{tab:sat-configurations}. BKS counts cases matching the best-known cycle time; ties are included. Proven counts solver-reported optimal solutions. Mean RPD uses runs with a feasible solution; PAR-2 uses all 72 cases. Bold values are best within a power limit; a metric tied across all six configurations is not emphasized.
\end{tablenotes}
\end{threeparttable}
\end{table}


Every two-phase configuration returns an incumbent in all 72 cases. Under
\(\Qmax^{\mathrm{top}}\), Baseline proves 44 cases, while SM and SM-T prove
43 each. SM-T reaches the BKS in 71 cases, and SM has the shortest median
times to the best incumbent and to proof. Under
\(\Qmax^{\mathrm{avg}}\), SM-T proves 24 cases and reaches the BKS in 64,
compared with 23 proofs and 61--62 BKS hits for the other encodings. It also
has the shortest median time to the best incumbent, whereas SM has the
shortest median proof time. Thus, none of the three encodings dominates every
solution-quality and proof metric.

\input{generated_table_estar_ablation}

Table~\ref{tab:estar-ablation} summarizes the three encoding pairs without
ranking the encodings themselves. Under \(\Qmax^{\mathrm{top}}\), \(E^*\)
increases the median clause count by 20.57--35.68\%; proof coverage is
unchanged for two encodings and lower for one, while the geometric-mean
runtime favors \(E\) for all three. Under \(\Qmax^{\mathrm{avg}}\), the
clause-count increase is 20.98--37.37\%; \(E^*\) improves proof coverage for
one encoding and runtime for two. Its benefit is therefore cap-dependent and
comes with consistently larger formulas. We retain the direct arc set \(E\)
in the locked two-phase campaign.

\subsection{RQ3: SAT versus MIP and CP}

% This table is generated by analysis/analyze_manuscript_results.py.
% Do not edit numerical values manually.
\begin{table}[t!]
\centering
\caption{Full-set performance of the selected SAT configuration and the commercial MIP and CP comparators.}
\label{tab:solver-overall}
\footnotesize
\renewcommand{\arraystretch}{1.08}
\setlength{\tabcolsep}{4pt}
\begin{threeparttable}
\begin{tabular*}{\textwidth}{@{\extracolsep{\fill}} l c c c c c @{}}
\toprule
Method & Feasible & BKS & Proven & Mean RPD & PAR-2 (s) \\
\midrule
\multicolumn{6}{@{}l}{\textit{Top-$m$-based threshold ($Q^{\mathrm{top}}_{\max}$)}} \\
\addlinespace[2pt]
Direct SAT (SM/E) & \textbf{72 / 72} & \textbf{71 / 72} & \textbf{44 / 72} & \textbf{0.01\%} & \textbf{2911.2} \\
Gurobi-MIP & 44 / 72 & 33 / 72 & 32 / 72 & 1.01\% & 4064.4 \\
CPLEX-MIP & 29 / 72 & 25 / 72 & 25 / 72 & 1.22\% & 4751.8 \\
CPLEX-CP & 16 / 72 & 16 / 72 & 16 / 72 & -- & 5603.1 \\
\addlinespace[4pt]
\multicolumn{6}{@{}l}{\textit{Average-based threshold ($Q^{\mathrm{avg}}_{\max}$)}} \\
\addlinespace[2pt]
Two-phase SAT (SM-T/E) & \textbf{72 / 72} & \textbf{64 / 72} & \textbf{24 / 72} & \textbf{0.11\%} & \textbf{4903.1} \\
Gurobi-MIP & 44 / 72 & 23 / 72 & 20 / 72 & 1.48\% & 5208.7 \\
CPLEX-MIP & 33 / 72 & 24 / 72 & 20 / 72 & 3.35\% & 5218.6 \\
CPLEX-CP & 14 / 72 & 14 / 72 & 14 / 72 & -- & 5800.1 \\
\bottomrule
\end{tabular*}
\begin{tablenotes}[flushleft]
\scriptsize
\item Feasible, BKS, and Proven respectively count recorded feasible solutions, best-known-solution hits, and proven optima out of 72 cases. For CPLEX-CP, an objective is archived only for the 16 top-$m$ and 14 average-based runs returning \textsc{optimal}; the other 56 and 58 runs, respectively, retain no incumbent. Its Feasible and BKS entries are therefore conservative recorded counts, and mean RPD is not reported. PAR-2 assigns twice the time limit to a run that does not prove optimality. Bold values are best among comparable methods within a power limit; an all-way tie is not emphasized.
\end{tablenotes}
\end{threeparttable}
\end{table}


Baseline/E is the selected two-phase SAT configuration under
\(\Qmax^{\mathrm{top}}\): it finds a feasible solution in all 72 cases and
proves 44, compared with 32 for Gurobi-MIP, 25 for CPLEX-MIP and 16 for
CPLEX-CP. Under \(\Qmax^{\mathrm{avg}}\), SM-T/E proves 24 cases, compared
with 20, 20 and 14 for the three baselines.

% This table is generated by analysis/analyze_manuscript_results.py.
% Do not edit numerical values manually.
\begin{table}[t!]
\centering
\caption{Paired comparison on cases for which SAT and a commercial comparator prove the same optimum.}
\label{tab:solver-paired}
\footnotesize
\renewcommand{\arraystretch}{1.08}
\setlength{\tabcolsep}{4pt}
\begin{threeparttable}
\begin{tabular*}{\textwidth}{@{\extracolsep{\fill}} l c c c c c @{}}
\toprule
Comparator & Common & SAT W/L & SAT-only & Comp.-only & Ratio ($t_{\mathrm{Comp}}/t_{\mathrm{SAT}}$) \\
\midrule
\multicolumn{6}{@{}l}{\textit{Top-$m$-based threshold ($Q^{\mathrm{top}}_{\max}$)}} \\
\addlinespace[2pt]
Gurobi-MIP & 31 & \textbf{31 / 0} & \textbf{13} & 1 & \textbf{27.81} \\
CPLEX-MIP & 25 & \textbf{25 / 0} & \textbf{19} & 0 & \textbf{68.73} \\
CPLEX-CP & 16 & \textbf{16 / 0} & \textbf{28} & 0 & \textbf{27.21} \\
\addlinespace[4pt]
\multicolumn{6}{@{}l}{\textit{Average-based threshold ($Q^{\mathrm{avg}}_{\max}$)}} \\
\addlinespace[2pt]
Gurobi-MIP & 19 & \textbf{18 / 1} & \textbf{5} & 1 & \textbf{14.25} \\
CPLEX-MIP & 19 & \textbf{18 / 1} & \textbf{5} & 1 & \textbf{30.12} \\
CPLEX-CP & 14 & \textbf{14 / 0} & \textbf{10} & 0 & \textbf{19.04} \\
\bottomrule
\end{tabular*}
\begin{tablenotes}[flushleft]
\scriptsize
\item Common counts cases with the same proven optimum. W/L counts SAT time wins / losses; Ratio is the geometric mean of $t_{\mathrm{Comp}}/t_{\mathrm{SAT}}$, so a value above one favors SAT. SAT-only and Comp.-only count cases proved only by the indicated method. Bold identifies a paired outcome favoring SAT: more time wins than losses, more SAT-only than Comp.-only proofs, or a ratio above one.
\end{tablenotes}
\end{threeparttable}
\end{table}


On common cases with the same proven optimum under
\(\Qmax^{\mathrm{top}}\), SAT wins every timed comparison; the geometric-mean
baseline/SAT ratios range from 24.07 to 61.24. Under
\(\Qmax^{\mathrm{avg}}\), the ratios range from 14.25 to 30.12, with one loss
against each MIP solver on the common proved set.

% This figure is generated by analysis/analyze_manuscript_results.py.
% Do not edit numerical coordinates manually.
\begin{figure}[H]
\centering
\begin{minipage}{0.49\textwidth}
\centering
\begin{tikzpicture}
\begin{axis}[width=\linewidth,height=45mm,xmode=log,xmin=0.003,xmax=3600,ymin=0,ymax=48,ytick={0,12,24,36,48},xlabel={Runtime (s)},ylabel={Proven cases},grid=major,grid style={gray!22,densely dotted},legend style={font=\scriptsize,at={(0.5,1.04)},anchor=south,legend columns=2,draw=gray!40,fill=white,fill opacity=0.95},tick label style={font=\scriptsize},label style={font=\small}]
\addplot[const plot,no marks,color=PlotBlue,solid,line width=1.3pt] coordinates {(0.003,0) (0.003541470,1) (0.005247116,2) (0.005295753,3) (0.008251190,4) (0.009734154,5) (0.010025740,6) (0.012866020,7) (0.012958288,8) (0.013652802,9) (0.016139984,10) (0.018746376,11) (0.025910139,12) (0.028381348,13) (0.034771204,14) (0.130478382,15) (0.204934835,16) (0.214081764,17) (0.278077841,18) (0.403787374,19) (0.546904325,20) (0.643606663,21) (0.673460960,22) (0.963696957,23) (1.005735874,24) (1.652459383,25) (2.954896927,26) (3.520453930,27) (5.610790014,28) (10.350619320,29) (14.979071380,30) (25.892252450,31) (30.701663730,32) (31.533023600,33) (45.685124640,34) (62.390542270,35) (69.248304370,36) (73.714394090,37) (102.675891400,38) (105.555170500,39) (427.714209100,40) (692.269913900,41) (1415.457757000,42) (1829.345566000,43) (3046.984778000,44)};
\addlegendentry{Direct SAT (SM/E)}
\addplot[const plot,no marks,color=PlotOrange,densely dashed,line width=1.1pt] coordinates {(0.003,0) (0.015736103,1) (0.038957596,2) (0.054211378,3) (0.094443798,4) (0.098404884,5) (0.104360580,6) (0.117602825,7) (0.152279139,8) (0.157175779,9) (0.260886192,10) (0.588612795,11) (0.619363069,12) (0.637070417,13) (0.646924496,14) (4.112289906,15) (4.747301102,16) (5.394608736,17) (19.196970460,18) (27.183960200,19) (29.504277940,20) (31.556227680,21) (76.154961110,22) (105.969540800,23) (185.230606300,24) (248.316510400,25) (258.247517600,26) (338.200370300,27) (456.894193200,28) (514.850428300,29) (637.906011300,30) (762.625196500,31) (930.133920000,32)};
\addlegendentry{Gurobi-MIP}
\addplot[const plot,no marks,color=PlotGreen,dashdotted,line width=1.1pt] coordinates {(0.003,0) (0.103470325,1) (0.113736629,2) (0.145992041,3) (0.160179615,4) (0.174883127,5) (0.249468327,6) (0.331637383,7) (0.336660147,8) (0.371800423,9) (0.509796143,10) (0.684638023,11) (0.715104818,12) (0.908212662,13) (1.468562365,14) (15.349610570,15) (26.812348370,16) (70.230240350,17) (115.213775600,18) (125.066116800,19) (146.385148800,20) (239.075233700,21) (267.683440700,22) (325.308251400,23) (435.926925400,24) (1953.797598000,25)};
\addlegendentry{CPLEX-MIP}
\addplot[const plot,no marks,color=PlotPurple,densely dotted,line width=1.1pt] coordinates {(0.003,0) (0.085296631,1) (0.101630449,2) (0.104367733,3) (0.120749712,4) (0.169259787,5) (0.202065468,6) (0.206754208,7) (0.216585875,8) (0.344884157,9) (0.394920349,10) (0.418884277,11) (0.458345175,12) (0.567824841,13) (1.174881220,14) (33.635968920,15) (185.564700400,16)};
\addlegendentry{CPLEX-CP}
\end{axis}
\end{tikzpicture}
\par\smallskip\footnotesize (a) $Q^{\mathrm{top}}_{\max}$
\end{minipage}
\hfill
\begin{minipage}{0.49\textwidth}
\centering
\begin{tikzpicture}
\begin{axis}[width=\linewidth,height=45mm,xmode=log,xmin=0.003,xmax=3600,ymin=0,ymax=48,ytick={0,12,24,36,48},xlabel={Runtime (s)},grid=major,grid style={gray!22,densely dotted},legend style={font=\scriptsize,at={(0.5,1.04)},anchor=south,legend columns=2,draw=gray!40,fill=white,fill opacity=0.95},tick label style={font=\scriptsize},label style={font=\small}]
\addplot[const plot,no marks,color=PlotBlue,solid,line width=1.3pt] coordinates {(0.003,0) (0.005237341,1) (0.006369829,2) (0.007310152,3) (0.009406328,4) (0.010154724,5) (0.011291504,6) (0.018010616,7) (0.018234015,8) (0.019078970,9) (0.021774769,10) (0.026400328,11) (0.046944141,12) (0.066911459,13) (0.067498446,14) (0.213874578,15) (0.268489361,16) (0.272375822,17) (1.245343924,18) (5.948643684,19) (22.627681020,20) (22.932389020,21) (586.137917000,22) (3200.426179000,23) (3586.001296000,24)};
\addlegendentry{2P-SAT (SM-T/E)}
\addplot[const plot,no marks,color=PlotOrange,densely dashed,line width=1.1pt] coordinates {(0.003,0) (0.014233112,1) (0.029942036,2) (0.049777985,3) (0.077031136,4) (0.108666658,5) (0.113059759,6) (0.213433027,7) (0.263217449,8) (0.323802948,9) (0.370296955,10) (0.388314009,11) (1.121779442,12) (1.308657169,13) (3.058547735,14) (7.559943438,15) (13.677929160,16) (23.913820980,17) (52.798798320,18) (74.101635220,19) (443.977200000,20)};
\addlegendentry{Gurobi-MIP}
\addplot[const plot,no marks,color=PlotGreen,dashdotted,line width=1.1pt] coordinates {(0.003,0) (0.074887753,1) (0.145617247,2) (0.147617817,3) (0.161184788,4) (0.302835226,5) (0.374069214,6) (0.389798164,7) (0.411795378,8) (0.539725065,9) (0.552820683,10) (0.633949518,11) (1.739472628,12) (3.550043344,13) (3.839728355,14) (27.962858680,15) (40.030714270,16) (274.104711100,17) (280.538741800,18) (290.674826600,19) (414.209457200,20)};
\addlegendentry{CPLEX-MIP}
\addplot[const plot,no marks,color=PlotPurple,densely dotted,line width=1.1pt] coordinates {(0.003,0) (0.084149361,1) (0.132140398,2) (0.182720661,3) (0.184361219,4) (0.211303949,5) (0.212219954,6) (0.225722313,7) (0.250324011,8) (0.311642885,9) (0.468177557,10) (0.575646639,11) (0.937197447,12) (1.363550663,13) (1.621778727,14)};
\addlegendentry{CPLEX-CP}
\end{axis}
\end{tikzpicture}
\par\smallskip\footnotesize (b) $Q^{\mathrm{avg}}_{\max}$
\end{minipage}
\caption{Cumulative number of cases solved to proven optimality within the 3,600-second time limit. The top-$m$-based panel uses Direct SAT (SM/E), whereas the average-based panel uses 2P-SAT (SM-T/E). Only runs returning \textsc{optimal} contribute; staircase height gives the number of cases proved by the corresponding wall-clock time. Both panels use the same axes.}
\label{fig:cactus}
\end{figure}


The cactus plots confirm that the coverage differences are not caused by a
few isolated fast cases. The selected SAT encodings reach higher final
staircases under both caps and generally reach common proof counts earlier.
Across the two regimes, \HardCaseCount{} case--cap combinations are proved by
at least one primary SAT encoding on \(E\) while none of the three MIP/CP
baselines proves optimality.

\subsection{RQ4: where do proofs become difficult?}

% This table is generated by analysis/analyze_manuscript_results.py.
% Do not edit numerical values manually.
\begin{table}[t!]
\centering
\caption{Cases solved to proven optimality by the selected SAT configuration, grouped by instance characteristics. Entries are proved / total.}
\label{tab:hardness-strata}
\small
\renewcommand{\arraystretch}{1.15}
\setlength{\tabcolsep}{28pt}
\begin{tabular}{@{} l c c @{}}
\toprule
Instance stratum & $Q^{\mathrm{top}}_{\max}$ & $Q^{\mathrm{avg}}_{\max}$ \\
\midrule
\multicolumn{3}{@{}l}{\textit{Tasks per station ($n/m$)}} \\
\addlinespace[2pt]
\quad $n/m\leq2$ & 8 / 10 & 7 / 10 \\
\quad $2<n/m\leq4$ & 23 / 48 & 11 / 48 \\
\quad $n/m>4$ & 13 / 14 & 6 / 14 \\
\addlinespace[4pt]
\multicolumn{3}{@{}l}{\textit{Direct precedence density ($d_E$)}} \\
\addlinespace[2pt]
\quad $d_E\leq0.10$ & 16 / 44 & 0 / 44 \\
\quad $0.10<d_E\leq0.20$ & 11 / 11 & 7 / 11 \\
\quad $d_E>0.20$ & 17 / 17 & 17 / 17 \\
\addlinespace[4pt]
\multicolumn{3}{@{}l}{\textit{Normalized energy lower-bound ratio ($\rho_E$)}} \\
\addlinespace[2pt]
\quad $\rho_E<1$ & 14 / 14 & 1 / 1 \\
\quad $1\leq\rho_E\leq1.25$ & 29 / 46 & 10 / 10 \\
\quad $\rho_E>1.25$ & 1 / 12 & 13 / 61 \\
\bottomrule
\end{tabular}
\end{table}


Table~\ref{tab:hardness-strata} is descriptive rather than causal. The
clearest separation is precedence density: under
\(\Qmax^{\mathrm{avg}}\), the selected SAT encoding proves none of the 44
sparsest cases but all 17 cases with direct-edge density above 0.20. Energy
pressure also distinguishes the regimes. Under the average-based cap, 61
cases have an energy lower bound more than 1.25 times the uncapped optimum,
and 13 are proved. This pattern is consistent with tight power budgets
making temporal coordination harder, while denser precedence relations prune
more station and timing combinations.

\subsection{Status of unproven best-known solutions}

All No-Peak references are independently certified, but proof coverage for
the capped problem is 45/72 under \(\Qmax^{\mathrm{top}}\) and 25/72 under
\(\Qmax^{\mathrm{avg}}\). The remaining BKS values are solver-reported
feasible UBs and are treated as such in every aggregate. In particular, AVG BUXEY-\(m=13\) has
BKS 42, AVG LUTZ2-\(m=49\) has \(C_0=20\) and BKS 19, and the corrected
SAWYER-\(m=12\) values discussed above are all unproven. No reported
\textsc{optimal} values disagree across solvers, no capped BKS is below its
exact No-Peak optimum, and the archived cap files reproduce their internal
minima. The new BKS 19 for AVG LUTZ2-\(m=49\) is reported by all three
two-phase encodings, but all three runs terminate at the time limit; it is
therefore retained as a feasible UB, not an optimum.
